% !TEX root = ../main.tex
\chapter{リファクタリングを等式として扱う}
\label{chap:ch34-refactoring-as-equations}

\begin{chapterabstract}
本章では、リファクタリングを「意味を変えない書き換え」として扱う。名前を変えた、関数を切り出した、共通関数に寄せた、処理順序を変えた、失敗を含む処理の括弧を付け替えた。このような変更は、外から見た振る舞いが同じであるなら、Lean では等式として表せる。本章の目的は、すべての実装をLeanへ移すことではない。変更で壊れてほしくない小さな意味保存性を選び、変更前後の関数が同じ結果を返すことを検査可能な形にすることである。
\end{chapterabstract}

\begin{learninggoals}
\item リファクタリングを「名前の変更」ではなく「観測可能な意味の保存」として説明できる。
\item 関数抽出、処理順序の変更、共通化、wrapper削除を等式として読み替えられる。
\item \texttt{Option} のような失敗を含む処理について、\texttt{bind} の括弧付け変更を意味保存として理解できる。
\item Lean 4 で小さな変更前後の関数を定義し、\texttt{rfl}、\texttt{rw}、\texttt{calc}、\texttt{cases} を使って安全性を証明できる。
\item 実務コードとLeanモデルの差分を明示し、コードレビューで使える証明の粒度を判断できる。
\end{learninggoals}

\section{この章を読む理由}

リファクタリングは、実装の内部構造を変える作業である。関数を分ける。重複した計算を共通関数へ寄せる。処理をパイプラインにする。不要になったadapterを削除する。失敗を含む処理を読みやすい \texttt{bind} の連鎖へ整える。

しかし、実務で本当に確認したいのは「コードの形がきれいになったか」だけではない。確認したい中心は、「外から見た意味が変わっていないか」である。入力が同じなら出力が同じか。成功・失敗の分岐は同じか。保存すべきID、件数、金額、状態、不変条件は変わっていないか。

\begin{intuitionbox}
リファクタリングを等式として見るとは、変更前の処理を左辺、変更後の処理を右辺に置き、両者が同じ結果を返すことを示す、ということである。つまり、コードレビューでよく言う「意味は変えていません」を、Lean が検査できる命題に変える。
\end{intuitionbox}

\FigurePlaceholder{fig-ch34-refactoring-equation-map.png}{0.90}{リファクタリングを等式として読む地図}{fig:ch34-refactoring-equation-map}

前の章までに、合成可能な設計、等式推論、関数外延性、モナドの直感を見てきた。本章ではそれらを実務のリファクタリングへ戻す。圏論の言葉でいえば、処理は射であり、リファクタリングは射の表現を変える作業である。変えるべきでないものは、射としての意味、つまり入力から出力への対応である。

\section{名前変更ではなく意味保存を見る}

まず、リファクタリングで保存したい性質を明確にする。関数名、変数名、ファイル分割、ディレクトリ構成は、直接には意味そのものではない。もちろん名前は読みやすさに影響するが、Leanで証明したいのは、名前の良し悪しではなく、観測できる振る舞いである。

\begin{definitionbox}
本章では、二つの処理 \(f\) と \(g\) が意味保存的に同じであるとは、扱うモデルの範囲で、すべての入力 \(x\) について \(f(x) = g(x)\) が成り立つことをいう。状態や失敗を含む処理では、返される状態、成功値、失敗値まで含めて等しいことを確認する。
\end{definitionbox}

たとえば、税込金額を計算する処理を考える。変更前は式を直接書いていた。変更後は、基礎金額の計算と税額計算を補助関数へ切り出した。この変更が安全であるとは、任意の小計と送料に対して、変更前後の合計が同じであるという意味である。

\begin{listing}[htbp]
\begin{minted}{lean4}
def basePrice (subtotal shipping : Nat) : Nat :=
  subtotal + shipping

def tax10 (n : Nat) : Nat :=
  n / 10

def totalBefore (subtotal shipping : Nat) : Nat :=
  (subtotal + shipping) + ((subtotal + shipping) / 10)

def totalAfter (subtotal shipping : Nat) : Nat :=
  let base := basePrice subtotal shipping
  base + tax10 base

theorem extract_base_price_preserves (subtotal shipping : Nat) :
    totalAfter subtotal shipping = totalBefore subtotal shipping := by
  rfl
\end{minted}
\caption{関数抽出による意味保存}
\label{lst:ch34-extract-base-price}
\end{listing}

\begin{leanbox}
この証明で \texttt{rfl} が通るのは、\texttt{totalAfter} を定義展開すると \texttt{totalBefore} と同じ式になるからである。ここでは高度な定理を使っていない。Leanに「補助関数を展開すれば同じ形になる」と確認させているだけである。
\end{leanbox}

このような証明は小さすぎるように見えるかもしれない。しかし、実務上の価値は「リファクタリングの単位を小さく切り出す」ことにある。複雑な関数全体ではなく、壊れると困る計算の核だけをLeanへ移す。そこで等式が通れば、少なくともその核については、変更が意味を変えていないと説明できる。

\section{関数抽出は補助定義の導入である}

関数抽出は、リファクタリングの中でも最も扱いやすい。多くの場合、変更前に式の中へ直接書かれていた部分式に名前を付けるだけだからである。数学的には、共通する部分式を補助定義にする操作と読める。

\begin{softwarebox}
実務で関数抽出を証明したいときは、まず「抽出した関数が何を計算しているか」を短い仕様にする。次に、抽出前の式と抽出後の式を並べ、補助関数を展開すれば同じ形になるかを確認する。
\end{softwarebox}

関数抽出が安全になりやすいのは、抽出した関数が純粋な関数である場合である。つまり、同じ入力なら同じ出力を返し、外部状態、時刻、乱数、DB、ネットワーク、ログ出力などに依存しない場合である。この条件が崩れると、同じ式に見えても意味が変わる可能性がある。

\begin{warningbox}
「同じコードを関数へ切り出しただけ」に見えても、評価タイミングが変わると意味が変わることがある。たとえば、現在時刻を読む処理、乱数を生成する処理、DBから値を読む処理、副作用を持つログ出力は、単純な等式として扱えないことがある。その場合は、時刻や外部入力を明示的な引数にするなど、証明しやすいモデルへ切り出す必要がある。
\end{warningbox}

この注意は、圏論の言葉では「同じ射として扱える範囲を決める」と言える。副作用を隠したままでは、射の入力と出力が足りない。Leanモデルでは、観測したい外部情報を型に入れ、処理を明示的な関数として表す。

\section{処理順序の変更は交換法則を要求する}

次に、処理順序の変更を考える。パイプラインの順序を入れ替えるリファクタリングは、見た目以上に危険である。二つの処理 \(f\) と \(g\) について、\texttt{g (f x)} と \texttt{f (g x)} が常に等しいとは限らないからである。

安全に順序を入れ替えられるのは、二つの処理が互いに干渉せず、等式
\[
  g(f(x)) = f(g(x))
\]
を満たす場合である。この形は、自然性や可換図式の章で見た「順序を変えても結果が同じ」という考え方と同じである。

\FigurePlaceholder{fig-ch34-order-change.png}{0.88}{処理順序の入れ替えと可換性}{fig:ch34-order-change}

単純な例として、金額に二種類の固定手数料を加える処理を考える。どちらも単に自然数を足すだけなら、順序を入れ替えても結果は同じである。ここで使う道具は、足し算の結合律と交換律である。

\begin{listing}[htbp]
\begin{minted}{lean4}
def addServiceFee (fee amount : Nat) : Nat :=
  amount + fee

def addHandlingFee (fee amount : Nat) : Nat :=
  amount + fee

theorem independent_fee_order (amount service handling : Nat) :
    addHandlingFee handling (addServiceFee service amount) =
    addServiceFee service (addHandlingFee handling amount) := by
  unfold addHandlingFee addServiceFee
  calc
    (amount + service) + handling = amount + (service + handling) := by
      rw [Nat.add_assoc]
    _ = amount + (handling + service) := by
      rw [Nat.add_comm service handling]
    _ = (amount + handling) + service := by
      rw [← Nat.add_assoc]
\end{minted}
\caption{独立した二つの手数料処理は順序を入れ替えられる}
\label{lst:ch34-order-preserves}
\end{listing}

\begin{leanbox}
この証明では、\texttt{unfold} で二つの補助関数を展開し、\texttt{calc} で式変形を段階的に示している。\texttt{Nat.add\_assoc} は結合律、\texttt{Nat.add\_comm} は交換律である。定理は暗記するものではなく、「この順序変更が安全である理由」を表す道具として使う。
\end{leanbox}

一方で、順序を入れ替えてはいけない処理も多い。たとえば「1を足してから2倍する」と「2倍してから1を足す」は一致しない。

\begin{listing}[htbp]
\begin{minted}{lean4}
def double (n : Nat) : Nat := n * 2

def addOne (n : Nat) : Nat := n + 1

example : double (addOne 3) ≠ addOne (double 3) := by
  decide
\end{minted}
\caption{順序を変えると意味が変わる小例}
\label{lst:ch34-order-not-always-safe}
\end{listing}

この小例は、処理順序変更のレビューで見るべき点をよく表している。すべての順序変更が危険なのではない。必要な法則がある順序変更だけが安全なのである。

\section{共通化は分岐の境界を明示する}

共通化もよくあるリファクタリングである。似たような関数が二つあり、重複している部分を一つの共通関数へ寄せる。このとき注意すべきなのは、共通化により「今まで別々だった仕様」が暗黙に混ざってしまうことである。

たとえば、クーポンなしの計算とクーポンありの計算を共通関数へ寄せるとする。クーポンなしの旧関数が、共通関数の \texttt{false} 分岐と一致することを示せれば、その分岐については意味保存が確認できる。

\begin{listing}[htbp]
\begin{minted}{lean4}
def discountCore (useCoupon : Bool) (price : Nat) : Nat :=
  if useCoupon then price - 10 else price

def oldWithoutCoupon (price : Nat) : Nat :=
  price

def newWithoutCoupon (price : Nat) : Nat :=
  discountCore false price

theorem commonize_without_coupon_preserves (price : Nat) :
    newWithoutCoupon price = oldWithoutCoupon price := by
  rfl
\end{minted}
\caption{共通化後の特定分岐が旧実装と一致することを示す}
\label{lst:ch34-commonize-preserves}
\end{listing}

この例では、\texttt{false} 分岐が旧実装と同じであることだけを示している。クーポンありの場合の仕様は別に検討する必要がある。共通化の証明では、すべてを一つの大きな定理に押し込めるより、分岐ごとに小さな定理へ分ける方がレビューしやすい。

\begin{softwarebox}
共通化のレビューでは、「共通関数が正しいか」だけでなく、「旧関数Aは共通関数のどの引数・どの分岐に対応するか」「旧関数Bはどの分岐に対応するか」を明示するとよい。Leanでは、それぞれを別々の定理にすることで、対応関係が読みやすくなる。
\end{softwarebox}

\section{不要なwrapperの削除は恒等処理の除去である}

adapterやwrapperを削除するリファクタリングも、等式として扱える。特に、そのwrapperが実質的に何もしない処理であるなら、圏論でいう恒等射、つまり恒等関数に近い。

\begin{listing}[htbp]
\begin{minted}{lean4}
def identityAdapter (x : Nat) : Nat := x

def wrappedNormalize (x : Nat) : Nat :=
  identityAdapter x

def directNormalize (x : Nat) : Nat :=
  x

theorem remove_identity_wrapper_preserves (x : Nat) :
    wrappedNormalize x = directNormalize x := by
  rfl
\end{minted}
\caption{恒等的なwrapperの削除}
\label{lst:ch34-remove-wrapper}
\end{listing}

ここでのポイントは、wrapperが本当に恒等的かどうかである。ログを出す、メトリクスを送る、例外を捕捉する、型変換時に正規化する、といった処理が入っているなら、それは「何もしない」処理ではない。削除してよいかどうかは、何を観測対象にするかによって変わる。

\begin{warningbox}
「戻り値が同じ」だけでは十分でないことがある。ログ、監査証跡、メトリクス、キャッシュ、DB更新などが仕様の一部なら、それらも出力や状態としてモデルに入れる必要がある。Leanモデルで観測しないものは、Leanの証明では保証されない。
\end{warningbox}

この注意は、チーム開発で特に重要である。証明は強いが、証明したモデルの範囲を超えては語れない。wrapper削除の安全性を説明するときは、「戻り値については等しい」「ログについてはこのモデルでは扱っていない」のように、保証範囲を明示する。

\section{モナド則による effectful code の整理}

失敗を含む処理では、普通の関数合成だけでは足りない。\texttt{Option}、\texttt{Except}、\texttt{Result} のような型を使うと、成功した場合だけ次の処理へ進み、失敗したらそこで止まる、という形になる。このような処理の合成には \texttt{bind} が現れる。

実務コードでは、次のようなリファクタリングが起きる。

\begin{itemize}
\item ネストした失敗処理を、平坦な \texttt{bind} の連鎖にする。
\item 途中の一時変数を消して、処理をパイプライン化する。
\item 不要な \texttt{pure} や \texttt{some} を取り除く。
\item 括弧の位置を変えて読みやすくする。
\end{itemize}

これらを安全に行うための考え方が、モナド則である。本章では一般のモナドを形式化せず、\texttt{Option} だけを使って、括弧の付け替えが意味を変えないことを見る。

\FigurePlaceholder{fig-ch34-option-bind-assoc.png}{0.88}{\texttt{Option.bind} の括弧付け変更}{fig:ch34-option-bind-assoc}

次の例では、正の値か確認し、偶数か確認し、最後に値を増やす。変更前は、最初の \texttt{bind} の中に次の \texttt{bind} がネストしている。変更後は、先に二つの検査を合成してから、最後の処理をつないでいる。

\begin{listing}[htbp]
\begin{minted}{lean4}
def parsePositive (n : Nat) : Option Nat :=
  if n = 0 then none else some n

def checkEven (n : Nat) : Option Nat :=
  if n % 2 = 0 then some n else none

def enrich (n : Nat) : Option Nat :=
  some (n + 1)

def flowBefore (n : Nat) : Option Nat :=
  Option.bind (parsePositive n) (fun x =>
    Option.bind (checkEven x) enrich)

def flowAfter (n : Nat) : Option Nat :=
  Option.bind (Option.bind (parsePositive n) checkEven) enrich
\end{minted}
\caption{\texttt{Option.bind} の結合的なリファクタリング}
\label{lst:ch34-option-assoc}
\end{listing}

この時点では、二つの定義が同じかどうかはまだ主張していない。次に、\texttt{parsePositive n} が \texttt{none} か \texttt{some x} かで場合分けし、さらに \texttt{checkEven x} も場合分けする。すべてのケースで同じ結果になることを確認する。

\begin{listing}[htbp]
\begin{minted}{lean4}
theorem option_assoc_refactor_preserves (n : Nat) :
    flowBefore n = flowAfter n := by
  unfold flowBefore flowAfter
  cases parsePositive n with
  | none => rfl
  | some x =>
      cases checkEven x with
      | none => rfl
      | some y => rfl
\end{minted}
\caption{失敗を含む処理の括弧変更が意味を保存すること}
\label{lst:ch34-option-assoc-proof}
\end{listing}

\begin{leanbox}
この証明は、\texttt{Option} の結合律を一般定理として呼び出しているわけではない。失敗する場合と成功する場合を展開し、どちらの書き方でも同じ結果になることをケースごとに確認している。初学者向けには、抽象的なモナド則よりも、このような具体的な場合分けから読む方が理解しやすい。
\end{leanbox}

\texttt{pure} に相当する \texttt{some} についても、小さな等式として確認できる。値を \texttt{some} に入れてすぐ \texttt{bind} で次の関数へ渡すことは、直接その関数を呼ぶことと同じである。また、\texttt{Option} 値を \texttt{some} へ渡すだけの \texttt{bind} は、元の値を変えない。

\begin{listing}[htbp]
\begin{minted}{lean4}
theorem option_bind_pure_left (x : Nat) (f : Nat -> Option Nat) :
    Option.bind (some x) f = f x := by
  rfl

theorem option_bind_pure_right (mx : Option Nat) :
    Option.bind mx some = mx := by
  cases mx with
  | none => rfl
  | some x => rfl
\end{minted}
\caption{\texttt{Option} における \texttt{pure} 相当の整理規則}
\label{lst:ch34-option-pure-laws}
\end{listing}

これらの等式は、effectful codeのリファクタリング規則として読める。不要な \texttt{some} を消す、\texttt{bind} の括弧を変える、失敗処理を読みやすく平坦化する。こうした変更は、法則が成り立つ範囲では意味を変えない。

\section{Lean 4で小さなリファクタリング証明を書く手順}

実務のリファクタリングをLeanで扱うときは、いきなり本番コード全体を移植しない。まず、意味保存を確認したい核だけを小さなモデルにする。

\begin{softwarebox}
おすすめの手順は次の通りである。
\begin{enumerate}
\item 変更前の処理を小さな関数として書く。
\item 変更後の処理を別の小さな関数として書く。
\item 保存したい性質を、戻り値の等式、関数の等式、状態の等式、または成功・失敗結果の等式として選ぶ。
\item まず \texttt{rfl} や \texttt{simp} で通る形を探す。
\item 通らない場合は、必要な補題を作り、\texttt{rw} や \texttt{calc} で段階的に示す。
\item 分岐を含む場合は、\texttt{cases} で全ケースを明示する。
\end{enumerate}
\end{softwarebox}

次の表は、よくあるリファクタリングとLeanでの見方を対応させたものである。

\small
\setlength{\tabcolsep}{3pt}
\begin{longtable}{@{}p{0.22\linewidth}p{0.36\linewidth}p{0.36\linewidth}@{}}
\toprule
リファクタリング & Leanでの表現 & 主な証明道具 \\
\midrule
関数抽出 & 補助関数を展開すると同じ式になる & \texttt{rfl}, \texttt{unfold}, \texttt{simp} \\
処理順序の変更 & \texttt{g (f x) = f (g x)} を示す & \texttt{calc}, \texttt{rw}, 結合律, 交換律 \\
共通化 & 旧関数が共通関数の特定分岐と一致する & \texttt{rfl}, \texttt{cases}, 分岐ごとの補題 \\
wrapper削除 & 恒等的な処理を消す & \texttt{rfl}, 関数外延性 \\
失敗処理の整理 & \texttt{bind} の括弧や \texttt{pure} を整理する & \texttt{cases}, モナド則 \\
\bottomrule
\end{longtable}

ここで重要なのは、証明対象を「変更全体」ではなく「保存したい意味」にすることである。たとえば、巨大な注文処理全体をLeanで再実装する必要はない。合計金額の計算だけ、ID保存だけ、失敗時に状態が変わらないことだけを切り出してもよい。

\section{実務へ持ち込むときの境界線}

Leanで証明したことは、モデルの範囲では強い保証になる。しかし、モデルと本番実装の間には必ず差分がある。そこを曖昧にすると、証明の価値も危うくなる。

\begin{warningbox}
「Leanで証明した」と言うときは、何を証明したのかを明記する。たとえば、「自然数モデル上の合計金額について、関数抽出前後の戻り値が等しい」「\texttt{Option} でモデル化した失敗処理について、\texttt{bind} の括弧付け変更が結果を変えない」のように書く。現実のDB、例外、並行性、時刻、I/O、丸め誤差まで自動的に保証されるわけではない。
\end{warningbox}

実務では、次の問いを設計レビューに入れるとよい。

\begin{itemize}
\item 変更前後で保存したい観測結果は何か。
\item その観測結果は、戻り値だけか、ログや状態も含むか。
\item Leanモデルでは、外部依存をどの引数として表すか。
\item 証明した性質と通常テストで見る性質をどう分担するか。
\item 証明が壊れたとき、それは仕様変更なのか、実装ミスなのか。
\end{itemize}

この章の範囲では、リファクタリングの安全性を完全形式手法として大規模に運用するところまでは踏み込まない。次章では、仕様、実装、テスト、証明をどのように接続し、どの性質をLeanで扱うべきかを整理する。

\section{本章のまとめ}

リファクタリングの安全性は、変更前後の意味保存として等式で表せる。

関数抽出は、補助定義を展開すると元の式へ戻ることとして確認できる。

処理順序の変更には、交換できる理由が必要である。足し算のように法則がある場合は、\texttt{calc} と \texttt{rw} で根拠を示せる。

共通化では、旧関数が共通関数のどの分岐に対応するかを明示する。

wrapper削除は、wrapperが本当に恒等的である範囲でだけ安全である。

\texttt{Option.bind} の括弧付け変更や \texttt{some} の整理は、effectful code のリファクタリング規則として読める。

Leanで証明するのは、本番コード全体ではなく、壊れると困る小さな意味保存性でよい。
